%% A realistic CEUR-WS workshop paper, shaped like the ceurart class's own
%% sample__1col.tex. Used to test venue-shift; not a real paper.
\documentclass{ceurart}

\usepackage{listings}

\begin{document}

\copyrightyear{2026}
\copyrightclause{Copyright for this paper by its authors.
  Use permitted under Creative Commons License Attribution 4.0 International (CC BY 4.0).}

\conference{AutoTrust'26: Workshop on Trust in Automation, June 07--11, 2026, Ulm, Germany}

\title{Trust Calibration in Highly Automated Driving}
\tnotemark[1]
\tnotetext[1]{This work was funded by the German Research Foundation.}

\author[1]{Mark Colley}[%
orcid=0000-0001-5207-5029,
email=mark.colley@uni-ulm.de,
url=https://www.uni-ulm.de/,
]
\cormark[1]
\address[1]{Institute of Media Informatics, Ulm University, James-Franck-Ring, Ulm, 89081, Germany}

\author[2]{Jane Doe}[%
email=jane.doe@example.edu,
]
\address[2]{Example University, Boston, USA}

\cortext[1]{Corresponding author.}

\begin{abstract}
  Automated vehicles must communicate their intent to build appropriate trust.
  We report a within-subjects study with 24 participants examining how takeover
  requests shape trust over repeated exposure.
\end{abstract}

\begin{keywords}
  automated driving \sep
  trust \sep
  takeover request \sep
  user study
\end{keywords}

\maketitle

\section{Introduction}
Trust in automated vehicles is neither uniformly good nor uniformly
bad, as \citet{lee2004trust} argue. We examine how it develops across
repeated takeovers.

\begin{figure}[t]
  \centering
  \includegraphics[width=\linewidth]{figures/apparatus}
  \caption{The driving simulator used in the study.}
  \label{fig:apparatus}
\end{figure}

\section{Method}
We ran a within-subjects experiment with 24 participants.

\begin{table}[t]
  \caption{Participant demographics.}
  \label{tab:demographics}
  \begin{tabular}{lr}
    \toprule
    Measure & Value \\
    \midrule
    Participants & 24 \\
    Mean age & 27.4 \\
    \bottomrule
  \end{tabular}
\end{table}

\section{Results}
Trust increased over exposure, $F(2,46) = 8.1$, $p < .001$.

\section{Conclusion}
Repeated exposure calibrates trust.

\section*{Declaration on Generative AI}
The author(s) have not employed any Generative AI tools.

\bibliography{references}

\end{document}
